\begin{abstract}
Tumors consist of diverse cell populations exhibiting both genetic and functional heterogeneity, which together influence disease progression and treatment response.
In cancer research, genotype-to-phenotype mapping aims to associate specific genetic alterations in individual cells or clones with functional traits such as transcriptional programs.
Follicular lymphoma, a common non-Hodgkin lymphoma, displays pronounced clonal diversity, yet its genetic and phenotypic heterogeneity have often been studied separately.
A previously developed probabilistic graphical model CaClust integrated multimodal data, including whole-exome, single-cell RNA, and B-cell receptor sequencing, to jointly infer clone genotypes and assign cells to clones, enabling single-cell genotyping.
Building on this foundation, this study proposes a deep learning framework based on variational autoencoders (VAEs) to capture more complex genotype-to-phenotype relationships.
This approach aims to enhance the resolution and interpretability of these links and reveal the functional consequences of somatic evolution in follicular lymphoma.

% Nowotwory składają się z różnorodnych populacji komórek, wykazujących zarówno genetyczną, jak i funkcjonalną heterogeniczność, które wspólnie wpływają na przebieg choroby oraz odpowiedź na leczenie.
% W badaniach nad nowotworami mapowanie genotyp–fenotyp ma na celu powiązanie określonych zmian genetycznych w poszczególnych komórkach lub klonach z cechami funkcjonalnymi, takimi jak programy transkrypcyjne.
% Chłoniak grudkowy, będący jednym z najczęstszych chłoniaków nieziarniczych, cechuje się wyraźną różnorodnością klonalną, jednak jego heterogeniczność genetyczna i fenotypowa były często analizowane oddzielnie.
% Wcześniej opracowany probabilistyczny model graficzny CaClust integrował dane multimodalne, w tym sekwencjonowanie całego egzomu, RNA pojedynczych komórek oraz receptorów komórek B, w celu wspólnego wnioskowania genotypów klonów i przypisywania komórek do odpowiednich klonów, umożliwiając genotypowanie na poziomie pojedynczych komórek.
% W oparciu o tę metodologię, niniejsze badanie proponuje zastosowanie głębokiego uczenia z wykorzystaniem wariacyjnych autoenkoderów (VAE), aby uchwycić bardziej złożone zależności między genotypem a fenotypem.
% Podejście to ma na celu zwiększenie rozdzielczości i interpretowalności tych powiązań oraz ujawnienie funkcjonalnych konsekwencji ewolucji somatycznej w chłoniaku grudkowym.
\end{abstract}+